\chapter{Homomorphismen, Kongruenzen und Quotienten}

\section{Strukturerhaltende Abbildungen}

\begin{definition}[Homomorphismus und Isomorphismus]
Eine Abbildung $f:S\to T$ heißt \emph{Homomorphismus}, wenn
$f(ab)=f(a)f(b)$ für alle $a,b\in S$ gilt.
Vorhandene neutrale Elemente müssen nicht erhalten werden.
Ein injektiver Homomorphismus heißt \emph{Einbettung}, ein
bijektiver Homomorphismus \emph{Isomorphismus}.
Die Schreibweise $S\cong T$ bedeutet, daß ein Isomorphismus existiert.
Ein Homomorphismus $S\to S$ heißt \emph{Endomorphismus}.
Für $B\subseteq T$ bezeichnet $f^{-1}(B)=\{s\in S:f(s)\in B\}$
das Urbild, auch wenn $f$ nicht bijektiv ist.
\end{definition}

\begin{theorem}[Komposition, Umkehrung und Produkte]
Kompositionen von Homomorphismen sind Homomorphismen.
Die Identitätsabbildung $\operatorname{id}_S$ ist ein Homomorphismus;
die Umkehrabbildung
eines Isomorphismus ist ein Homomorphismus.
Für jeden Homomorphismus $f:S\to T$ gilt
\[
 f(a_1\cdots a_n)=f(a_1)\cdots f(a_n)\qquad(n\geq1).
\]
Insbesondere ist $f(a^n)=f(a)^n$.
\cite{harju}
\end{theorem}

\begin{theorem}[Bilder, Urbilder und Bestimmung durch Erzeuger]
Bilder und Urbilder von Unterhalbgruppen unter einem Homomorphismus
sind Unterhalbgruppen. Für $f:S\to T$ und $X\subseteq S$ gilt $f(\langle X\rangle)=\langle f(X)\rangle$.

Zwei Homomorphismen $S\to T$, die auf einem Erzeugendensystem
übereinstimmen, stimmen auf ganz $S$ überein.
\cite{harju,gould}
\end{theorem}

\begin{theorem}[Universelle Eigenschaft des direkten Produkts]
Die Koordinatenprojektionen $\pi_i:\prod_{j\in I}S_j\to S_i$
sind Homomorphismen. Zu jeder Familie von Homomorphismen
$f_i:T\to S_i$ gibt es genau einen Homomorphismus
$f:T\to\prod_i S_i$ mit $\pi_i\circ f=f_i$ für alle $i$;
er ist durch $f(t)=(f_i(t))_i$ gegeben.
\cite{howie}
\end{theorem}

\begin{theorem}[Transport auf Potenzhalbgruppen]
Für einen Homomorphismus $f:S\to T$ ist
\[
 \mathcal P(f):\mathcal P(S)\to\mathcal P(T),\qquad
 A\mapsto f(A)
\]
ein Homomorphismus und erhält beliebige Vereinigungen.
Ist $f$ injektiv, surjektiv oder bijektiv, so gilt dies entsprechend
für $\mathcal P(f)$. Außerdem gelten
$\mathcal P(g\circ f)=\mathcal P(g)\circ\mathcal P(f)$ und
$\mathcal P(\mathrm{id}_S)=\mathrm{id}_{\mathcal P(S)}$.
Insbesondere folgt aus $S\cong T$ stets
$\mathcal P(S)\cong\mathcal P(T)$.
\cite{projectpower}
\end{theorem}

\section{Kongruenzen und der Homomorphiesatz}

\begin{definition}[Kongruenz und Kernrelation]
Eine Äquivalenzrelation $\rho$ auf $S$ heißt \emph{Kongruenz}, wenn
\[
 a\mathrel{\rho}b\ \Longrightarrow\
 xa\mathrel{\rho}xb\ \text{und}\ ax\mathrel{\rho}bx
 \qquad(a,b,x\in S).
\]
Die Menge aller Kongruenzen heißt $\operatorname{Con}(S)$.
Für einen Homomorphismus $f:S\to T$ heißt $\ker f=\{(a,b)\in S\times S:f(a)=f(b)\}$ seine \emph{Kernrelation}.

Es sei $\Delta_S=\{(a,a):a\in S\}$.
\end{definition}

\begin{theorem}[Quotientenbildung und universelle Eigenschaft]
\label{b28guide:thm:quotient}
Für jede Kongruenz $\rho$ ist die Vorschrift
\[
 [a]_\rho[b]_\rho=[ab]_\rho
\]
wohldefiniert und macht $S/\rho$ zu einer Halbgruppe.
Die Abbildung $q_\rho:S\to S/\rho$, $a\mapsto[a]_\rho$, ist
ein surjektiver Homomorphismus mit $\ker q_\rho=\rho$.
Ein Homomorphismus $f:S\to T$ besitzt genau dann eine
Faktorisierung $f=\bar f\circ q_\rho$ durch einen Homomorphismus
$\bar f:S/\rho\to T$, wenn $\rho\subseteq\ker f$.
In diesem Fall ist $\bar f$ eindeutig und durch
$\bar f([a]_\rho)=f(a)$ gegeben.
\cite{harju,gould}
\end{theorem}

\begin{theorem}[Erster Isomorphiesatz]
\label{b28guide:thm:first-iso}
Für jeden Homomorphismus $f:S\to T$ ist $\ker f$ eine Kongruenz,
und die Abbildung
\[
 S/\ker f\longrightarrow f(S),\qquad [a]_{\ker f}\longmapsto f(a)
\]
ist ein Isomorphismus. Jede homomorphe Abbildung zerfällt damit
in einen surjektiven Quotientenhomomorphismus und eine Einbettung.
\cite{harju,gould}
\end{theorem}

\begin{theorem}[Einschränkung auf Unterhalbgruppen]
Für eine Unterhalbgruppe $A\subseteq S$ und eine Kongruenz $\rho$
ist $\rho|_A=\rho\cap(A\times A)$ eine Kongruenz auf $A$, und
\[
 A/(\rho|_A)\ \cong\ q_\rho(A)\subseteq S/\rho,
 \qquad [a]_{\rho|_A}\longmapsto[a]_\rho .
\]
\cite{howie}
\end{theorem}

\begin{theorem}[Korrespondenzsatz und iterierte Quotienten]
\label{b28guide:thm:correspondence}
Für $\rho\in\operatorname{Con}(S)$ besteht ein Ordnungsisomorphismus
\[
 \{\sigma\in\operatorname{Con}(S):\rho\subseteq\sigma\}
 \longrightarrow\operatorname{Con}(S/\rho),\qquad
 \sigma\longmapsto\sigma/\rho ,
\]
wobei $[a]_\rho\mathrel{\sigma/\rho}[b]_\rho$ genau dann gilt,
wenn $a\mathrel{\sigma}b$.
Seine Umkehrung ist das Urbild unter $q_\rho\times q_\rho$.
Für $\rho\subseteq\sigma$ gilt $(S/\rho)/(\sigma/\rho)\cong S/\sigma$.

\cite{harju,gould}
\end{theorem}

\section{Erzeugte Kongruenzen}

\begin{definition}[Von einer Relation erzeugte Kongruenz]
Für $R\subseteq S\times S$ bezeichnet
$\operatorname{Cg}_S(R)$ den Durchschnitt aller Kongruenzen auf $S$,
die $R$ enthalten.
Die Schreibweise $R^{-1}=\{(v,u):(u,v)\in R\}$ bezeichnet
die umgekehrte Relation.
\end{definition}

\begin{theorem}[Vollständiger Kongruenzverband]
$\operatorname{Con}(S)$ ist unter Inklusion ein vollständiger Verband,
mit kleinstem Element $\Delta_S$, größtem Element $S\times S$ und
\[
 \bigwedge_i\rho_i=\bigcap_i\rho_i,\qquad
 \bigvee_i\rho_i=\operatorname{Cg}_S\left(\bigcup_i\rho_i\right).
\]
Der Durchschnitt einer leeren Kongruenzfamilie ist $S\times S$.
Insbesondere ist $\operatorname{Cg}_S(R)$ die kleinste Kongruenz,
die $R$ enthält.
\cite{harju,howie}
\end{theorem}

\clearpage
\begin{theorem}[Beschreibung durch endliche Kontextketten]
\label{b28guide:thm:chains}
Es gilt $a\mathrel{\operatorname{Cg}_S(R)}b$ genau dann, wenn
$a=b$ ist oder eine endliche Folge $a=z_0,z_1,\ldots,z_n=b$
existiert, deren aufeinanderfolgende Glieder jeweils die Form
\[
 z_i=xuy,\qquad z_{i+1}=xvy
\]
mit für jeden Schritt geeigneten $x,y\in S^1$ und
$(u,v)\in R\cup R^{-1}$ haben.
Jede einzelne Folgerung aus $R$ verwendet bereits eine endliche
Teilmenge von $R$:
\[
 \operatorname{Cg}_S(R)=
 \bigcup_{\substack{F\subseteq R\\F\text{ endlich}}}
 \operatorname{Cg}_S(F).
\]
Hier wird die Endlichkeit einer Ableitung festgestellt;
$S$ und $R$ dürfen beliebig unendlich sein.
\cite{harju}
\end{theorem}

\begin{theorem}[Algebraischer Kongruenzverband]
\label{b28guide:thm:algebraic-congruences}
Für jede Halbgruppe $S$ ist $\operatorname{Con}(S)$ algebraisch
im Sinne von Definition~\ref{b28guide:def:compact}.
Die kompakten Elemente sind genau die Kongruenzen
$\operatorname{Cg}_S(F)$ mit endlichem $F\subseteq S\times S$.
Jede Kongruenz ist die gerichtete Vereinigung der in ihr enthaltenen
endlich erzeugten Kongruenzen.

Für jede nichtleere, unter Inklusion gerichtete Familie
$(\rho_i)_{i\in I}$ gilt
\[
 \bigvee_{i\in I}\rho_i=\bigcup_{i\in I}\rho_i.
\]
Dabei bedeutet \emph{gerichtet}, daß zu $i,j\in I$ ein $k\in I$
mit $\rho_i\cup\rho_j\subseteq\rho_k$ existiert.
Eine Endlichkeit von $S$ wird nicht verlangt. \cite{burris}
\end{theorem}
